% !TEX root = IntegrationByParts.tex
%
% settings.json:
%
% "latex-workshop.latex.recipes": [
% { // Uses LaTeX
%     "name": "tex",
%     "tools": [
%         "lualatex"
%     ]
% },

\documentclass[10pt]{article}

% ---------------------------------------------------------------
% Page geometry
% ---------------------------------------------------------------
\usepackage[
  a4paper,
  top=0.82in,
  bottom=0.82in,
  left=0.95in,
  right=0.95in
]{geometry}

% ---------------------------------------------------------------
% Fonts and microtypography
% ---------------------------------------------------------------
\usepackage[T1]{fontenc}
\usepackage{lmodern}
\usepackage{microtype}

% ---------------------------------------------------------------
% Mathematics
% ---------------------------------------------------------------
\usepackage{amsmath,amssymb}
\newcommand{\dd}{\mathop{}\!\mathrm{d}}

% ---------------------------------------------------------------
% Spacing
% ---------------------------------------------------------------
\usepackage{setspace}
\setstretch{1.04}

\setlength{\parskip}{4pt}
\setlength{\parindent}{0pt}

% Tighter display math spacing (\normalsize resets these at
% \begin{document}, so patch it directly)
\makeatletter
\g@addto@macro\normalsize{%
  \setlength\abovedisplayskip{5pt}
  \setlength\belowdisplayskip{5pt}
  \setlength\abovedisplayshortskip{3pt}
  \setlength\belowdisplayshortskip{3pt}}
\makeatother

% ---------------------------------------------------------------
% Page style (no header, footer, or page number)
% ---------------------------------------------------------------
\pagestyle{empty}

% ---------------------------------------------------------------
% Date formatting
% ---------------------------------------------------------------
\renewcommand{\today}{%
  \number\day\space
  \ifcase\month\or
  January\or February\or March\or April\or May\or June\or July\or August\or September\or October\or November\or December
  \fi
  \space\number\year
}

% ---------------------------------------------------------------
\begin{document}

% ---------------------------------------------------------------
% Title block
% ---------------------------------------------------------------
\begin{center}
  {\large\textsc{A quick note on integration by parts}}\\[4pt]
  By \textsc{Hanson Char}\\
  {\small\textit{Revised \today}}
\end{center}

Recall the integration by parts formula:
\[
\int u \dd v = uv - \int v \dd u.
\]

Neat, but the use of this formula requires mapping \(u\) and \(v\) to the specific integrand. This can be cumbersome, especially when the integrand is a product of three or more functions. Now, consider a seemingly more complicated but equivalent formulation:
\[
\int xy \dd z = x \int y \dd z - \int \left( \int y \dd z \right) \dd x,
\]
which turns out to be easier to apply in practice. Why? Because it reduces the amount of mental book-keeping: the right-hand side can be written directly from the left-hand side. Specifically, given the pattern
\[
\int x \, (\text{something}) \dd z,
\]
we can rewrite it immediately as
\[
x \int (\text{something}) \dd z - \int \left( \int (\text{something}) \dd z \right) \dd x.
\]

Happy integrating by parts!

\begin{center}
---
\end{center}

\end{document}
